% BANKS-SOURCE: pdf=23 print=13 kind=solution-section id=chapter02-numbered-solutions
\begin{bankssolutions}{Solutions to Problems for Chapter 2}

% BANKS-SOURCE: pdf=23 print=13 kind=solution id=2.1
\BanksSolution{2.1}
Take $\eta=\operatorname{diag}(-1,1,\ldots,1)$. Taking determinants in
$\Lambda^{T}\eta\Lambda=\eta$ gives
\[
  (\det\Lambda)^{2}\det\eta=\det\eta,
  \qquad \det\Lambda=\pm1.
\]
The $00$ component of the same equation is
\[
  (\Lambda^{0}{}_{0})^{2}-\sum_{i=1}^{D-1}(\Lambda^{i}{}_{0})^{2}=1,
\]
so $|\Lambda^{0}{}_{0}|\geq1$. Hence neither $\det\Lambda$ nor the sign of $\Lambda^{0}{}_{0}$ can change along a continuous path in the Lorentz group. The identity has $(\det\Lambda,\operatorname{sgn}\Lambda^{0}{}_{0})=(+,+)$, which gives the proper orthochronous component.

Let
\[
  \mathsf T=\operatorname{diag}(-1,1,\ldots,1),
  \qquad
  \mathsf P_{1}=\mathbf 1_D-2\,\widehat e_{1}\widehat e_{1}^{T},
  \qquad \widehat e_{1}=(0,\mathbf e_{1})^{T}.
\]
Here $\mathbf e_{1}$ is the first spatial basis vector.
Both preserve $\eta$. Their determinants and time-component signs are
\[
  \det\mathsf T=-1,\quad \operatorname{sgn}\mathsf T^{0}{}_{0}=-,
  \qquad
  \det\mathsf P_{1}=-1,\quad \operatorname{sgn}\mathsf P_{1}^{0}{}_{0}=+,
\]
while $\mathsf T\mathsf P_{1}$ has $\det(\mathsf T\mathsf P_{1})=+1$ and time-component sign $-$. Left multiplication by these fixed transformations maps the proper orthochronous component onto the other sign sectors. Thus each sector is a connected component; the representatives are time reversal, a spatial reflection, and their product.

For completeness, the $(+,+)$ sector is connected. If $L$ lies in it, the future unit time-like vector $Le_{0}$ can be reached from $e_{0}$ by a boost $B$ with a continuously variable rapidity. The transformation $B^{-1}L$ fixes $e_{0}$ and is an element of $SO(D-1)$, which is connected. A path through $B$ followed by a path through this rotation connects $L$ to the identity.

For a nonzero time-like vector $v$, $v^{2}<0$, every Lorentz transform satisfies
\[
  ((\Lambda v)^{0})^{2}=|\boldsymbol{\Lambda v}|^{2}-v^{2}>0.
\]
For a nonzero null vector, the same equation gives $((\Lambda v)^{0})^{2}=|\boldsymbol{\Lambda v}|^{2}$; a zero time component would force the transformed vector to vanish. Along a path $\Lambda_s$ in the proper orthochronous component, $(\Lambda_s v)^{0}$ is continuous and never zero. Its sign therefore equals the sign of $v^{0}$ at the identity. This proves the time-orientation statement for time-like and null vectors.

For a space-like vector $q$, rotate its spatial part so that $q=(q^{0},q^{1},\mathbf 0_{D-2})$, with $|q^{1}|>|q^{0}|$. A boost in the $1$-direction can be written
\[
  B(\chi)=
  \begin{pmatrix}
    \cosh\chi&-\sinh\chi\\[-2pt]
    -\sinh\chi&\cosh\chi
  \end{pmatrix}
    \oplus\mathbf 1_{D-2},
  \qquad
  (B(\chi)q)^{0}=q^{0}\cosh\chi-q^{1}\sinh\chi.
\]
This is a proper orthochronous transformation. Since $|q^{0}/q^{1}|<1$, the choice
$\tanh\chi=q^{0}/q^{1}$ makes the time component zero. Values of $\chi$ on the two sides of this value give opposite signs, so a proper orthochronous transformation can also change the sign.

% BANKS-SOURCE: pdf=24 print=14 kind=solution id=2.2
\BanksSolution{2.2}
The $t_{0}$-dependent factor is constant under $\partial_{t}$. Differentiating the corrected definition and using the Schrödinger-picture equation gives
\begin{align*}
  \ii\partial_{t}U_{D}
  &= -H_{0}\ee^{\ii H_{0}t}U_{S}\ee^{-\ii H_{0}t_{0}}
     +\ee^{\ii H_{0}t}(H_{0}+V)U_{S}\ee^{-\ii H_{0}t_{0}}\\
  &=\ee^{\ii H_{0}t}V U_{S}\ee^{-\ii H_{0}t_{0}}.
\end{align*}
The corrected rightmost factor makes $U_D(t_{0},t_{0})=\mathbf 1$. Inserting the identity $\ee^{-\ii H_{0}t}\ee^{\ii H_{0}t}$ next to $V$ yields
\[
  \ii\partial_{t}U_{D}
  =\left(\ee^{\ii H_{0}t}V\ee^{-\ii H_{0}t}\right)U_{D}
  =V(t)U_{D}.
\]

% BANKS-SOURCE: pdf=24 print=14 kind=solution id=2.3
\BanksSolution{2.3}
The boundary condition is $U_{D}(t_{0},t_{0})=\mathbf 1$. Integrating the differential equation once gives
\[
  U_{D}(t,t_{0})=\mathbf 1-\ii\int_{t_{0}}^{t}\dd s_{1}\,V(s_{1})U_{D}(s_{1},t_{0}).
\]
Repeated substitution produces the Dyson series
\[
  U_{D}(t,t_{0})=\mathbf 1+
  \sum_{n=1}^{\infty}(-\ii)^{n}
  \int_{t_{0}}^{t}\dd s_{1}\int_{t_{0}}^{s_{1}}\dd s_{2}\cdots
  \int_{t_{0}}^{s_{n-1}}\dd s_{n}\,
  V(s_{1})V(s_{2})\cdots V(s_{n}).
\]
The integration region has the times ordered as $s_{1}\geq s_{2}\geq\cdots\geq s_{n}$. Replacing it by the full $n$-cube and inserting the time-ordering operator counts each ordered region $n!$ times. Therefore
\[
  U_{D}(t,t_{0})
  =\sum_{n=0}^{\infty}\frac{(-\ii)^{n}}{n!}
    \int_{t_{0}}^{t}\dd s_{1}\cdots\dd s_{n}\,
    T\!\left[V(s_{1})\cdots V(s_{n})\right]
  =T\ee^{-\ii\int_{t_{0}}^{t}V(s)\,\dd s}.
\]

% BANKS-SOURCE: pdf=24 print=14 kind=solution id=2.4
\BanksSolution{2.4}
For an oscillator with coordinate $q$, oscillator mass $M$, and frequency $\omega$, the normalized ground state is
\[
  \psi_{\omega}(q)=\left(\frac{M\omega}{\pi}\right)^{1/4}
  \exp\!\left(-\frac{M\omega q^{2}}{2}\right).
\]
The overlap for two frequencies is the Gaussian integral
\begin{align*}
  c(\omega_{1},\omega_{2})
  &=\int\dd q\,\psi_{\omega_{1}}(q)^{*}\psi_{\omega_{2}}(q)\\
  &=\left(\frac{2\sqrt{\omega_{1}\omega_{2}}}{\omega_{1}+\omega_{2}}\right)^{1/2}
   =\left[\cosh\!\left(\frac12\log\frac{\omega_{2}}{\omega_{1}}\right)\right]^{-1/2}.
\end{align*}
It is strictly smaller than one when the frequencies differ.

Put the scalar field in a spatial box of volume $\mathcal V$ in $D$ spacetime dimensions. Each independent real Fourier mode is an oscillator with
\[
  \omega_{\mathbf k,m}=\sqrt{|\mathbf k|^{2}+m^{2}}.
\]
Consequently, up to the harmless choice of independent representatives of the pairs $\mathbf k$ and $-\mathbf k$,
\[
  \langle0_{m_{1}}|0_{m_{2}}\rangle
  =\prod_{\mathbf k}c(\omega_{\mathbf k,m_{1}},\omega_{\mathbf k,m_{2}}),
  \qquad
  \log|\langle0_{m_{1}}|0_{m_{2}}\rangle|
  =\sum_{\mathbf k}\log c_{\mathbf k}.
\]
In infinite volume,
\[
  \sum_{\mathbf k}\longrightarrow
  \frac{\mathcal V}{(2\pi)^{D-1}}\int\dd^{D-1}\mathbf k.
\]
For $m_{1}\ne m_{2}$, the integrand is strictly negative on a set of nonzero measure. The logarithm therefore tends to $-\infty$ as $\mathcal V\to\infty$, with any fixed ultraviolet regulator, and the vacuum overlap is zero in every spatial dimension.

The ultraviolet behavior also gives a dimensional check on the finite-volume claim. For $|\mathbf k|\to\infty$,
\[
  \log c_{\mathbf k}
  =-\frac{(m_{2}^{2}-m_{1}^{2})^{2}}{64|\mathbf k|^{4}}
   +O(|\mathbf k|^{-6}).
\]
Thus the unregulated mode sum diverges in the ultraviolet for $D-1\geq4$ and converges for $D-1=1,2,3$. The problem's threshold has been corrected from the source's $d\geq2$ to $D-1\geq4$; for $D-1=2,3$ the displayed product has a nonzero ultraviolet limit at fixed volume.

% BANKS-SOURCE: pdf=24 print=14 kind=solution id=2.5
\BanksSolution{2.5}
Write $p^{0}=E_{\mathbf p}\equiv\omega_{\mathbf p}=\sqrt{\mathbf p^{2}+m^{2}}$, so that $p^{2}=-m^{2}$, and use the one-particle transformation law from the preceding discussion:
\[
  U(\Lambda)a^{\dagger}(\mathbf p)U(\Lambda)^{-1}
  =\sqrt{\frac{\omega_{\boldsymbol{\Lambda p}}}{\omega_{\mathbf p}}}
    a^{\dagger}(\boldsymbol{\Lambda p}),
  \qquad U(\Lambda)|0\rangle=|0\rangle.
\]
The annihilation-operator law is its adjoint. Since
\[
  \frac{\dd^{D-1}\mathbf p}{\omega_{\mathbf p}}=\frac{\dd^{D-1}\boldsymbol{\Lambda p}}{\omega_{\boldsymbol{\Lambda p}}},
  \qquad p\mathbin{\cdot}x=(\Lambda p)\mathbin{\cdot}(\Lambda x),
\]
changing variables from $p$ to $q=\Lambda p$ gives
\[
  U(\Lambda)\phi(x)U(\Lambda)^{-1}=\phi(\Lambda x).
\]
The square-root factor in the state normalization is exactly what converts the noncovariant $\dd^{D-1}\mathbf p$ measure into the invariant $\dd^{D-1}\mathbf p/\omega_{\mathbf p}$ measure.

Let $\varepsilon=+1$ for bosons and $\varepsilon=-1$ for fermions, and define the graded bracket by $[A,B]_{\varepsilon}=AB-\varepsilon BA$. With
\[
  \Delta^{+}(z)=\int\frac{\dd^{D-1}\mathbf p}{(2\pi)^{D-1}2\omega_{\mathbf p}}\,\ee^{\ii p\mathbin{\cdot}z},
\]
the field expansion gives
\[
  [\phi(x),\phi(y)]_{\varepsilon}
  =\Delta^{+}(x-y)-\varepsilon\Delta^{+}(y-x).
\]
For a space-like $z=x-y$, choose a frame with $z^{0}=0$. For bosons the integrand is odd under $\mathbf p\mapsto-\mathbf p$, so the commutator vanishes. For fermions the equal-time anticommutator contains the even sum of the two Fourier waves and is nonzero. The scalar field is therefore local with the ordinary commutator precisely for Bose statistics.

At equal time, direct differentiation gives
\begin{align*}
  [\phi(t,\mathbf x),\dot\phi(t,\mathbf y)]
  &=\int\frac{\dd^{D-1}\mathbf p}{(2\pi)^{D-1}2\omega_{\mathbf p}}
    \ii\omega_{\mathbf p}\left(\ee^{\ii\mathbf p\cdot(\mathbf x-\mathbf y)}
    +\ee^{-\ii\mathbf p\cdot(\mathbf x-\mathbf y)}\right)\\
  &=\ii\delta^{D-1}(\mathbf x-\mathbf y),
\end{align*}
and the same-time field commutator is zero. In $D$ spacetime dimensions $[\phi]=(D-2)/2$ and $[\dot\phi]=D/2$, so both sides of the second relation have mass dimension $D-1$. The field and its time derivative are canonical conjugates.

% BANKS-SOURCE: pdf=24 print=14 kind=solution id=2.6
\BanksSolution{2.6}
To test unequal masses, temporarily write the two frequencies separately:
\[
  \omega_{\pm,\mathbf p}=\sqrt{\mathbf p^{2}+m_{\pm}^{2}},
  \qquad p_{\pm}=(\omega_{\pm,\mathbf p},\mathbf p),\quad p_{\pm}^{2}=-m_{\pm}^{2}.
\]
The charged field and its adjoint then have the form
\[
  \phi(x)=\int\frac{\dd^{D-1}\mathbf p}{(2\pi)^{(D-1)/2}}
  \left(\frac{a(\mathbf p)\ee^{\ii p_{+}\mathbin{\cdot}x}}{\sqrt{2\omega_{+,\mathbf p}}}
  +\eta\frac{b^{\dagger}(\mathbf p)\ee^{-\ii p_{-}\mathbin{\cdot}x}}{\sqrt{2\omega_{-,\mathbf p}}}\right),
\]
\[
  \phi^{\dagger}(y)=\int\frac{\dd^{D-1}\mathbf q}{(2\pi)^{(D-1)/2}}
  \left(\frac{a^{\dagger}(\mathbf q)\ee^{-\ii q_{+}\mathbin{\cdot}y}}{\sqrt{2\omega_{+,\mathbf q}}}
  +\eta^{*}\frac{b(\mathbf q)\ee^{\ii q_{-}\mathbin{\cdot}y}}{\sqrt{2\omega_{-,\mathbf q}}}\right).
\]
Using independent particle and antiparticle oscillators, the only two nonzero terms in the commutator are
\[
  [\phi(x),\phi^{\dagger}(y)]
  =\Delta^{+}_{m_{+}}(x-y)-|\eta|^{2}\Delta^{+}_{m_{-}}(y-x).
\]
For a space-like separation choose equal time and Fourier transform in the spatial separation. Locality then requires, for every $\mathbf p$,
\[
  \frac{1}{2\omega_{+,\mathbf p}}
  =\frac{|\eta|^{2}}{2\omega_{-,\mathbf p}}.
\]
At large $|\mathbf p|$ both frequencies approach $|\mathbf p|$, so $|\eta|^{2}=1$. The equality of the functions of $|\mathbf p|$ then gives $m_{+}^{2}=m_{-}^{2}$; positive masses imply $m_{+}=m_{-}$. Hence $\eta=\ee^{\ii\vartheta}$ is a pure phase. Defining $b'=\eta^{*}b$ gives $b'^{\dagger}=\eta b^{\dagger}$ and removes the phase from $\phi$. The scalar locality computation also selects Bose statistics for the ordinary commutator.

% BANKS-SOURCE: pdf=25 print=15 kind=solution id=2.7
\BanksSolution{2.7}
With $V(t)=-\int\dd^{D-1}\mathbf x\,\mathcal L_{I}(t,\mathbf x)$, the infinite-time interaction-picture operator is
\[
  S=\lim_{t_{\pm}\to\pm\infty}
  T\exp\!\left(\ii\int_{t_{-}}^{t_{+}}\dd t\int\dd^{D-1}\mathbf x\,\mathcal L_{I}(t,\mathbf x)\right)
  =T\exp\!\left(\ii\int\dd^{D}x\,\mathcal L_{I}(x)\right).
\]
Its $n$th term is
\[
  S_{n}=\frac{\ii^{n}}{n!}\int\dd^{D}x_{1}\cdots\dd^{D}x_{n}\,
  T\!\left[\mathcal L_{I}(x_{1})\cdots\mathcal L_{I}(x_{n})\right].
\]
For a Lorentz scalar, $U(\Lambda)\mathcal L_{I}(x)U(\Lambda)^{-1}=\mathcal L_{I}(\Lambda x)$, and the $D$-volume measure is invariant. A proper orthochronous transformation preserves the order of two time-like separated points. It can change the coordinate-time order of space-like separated points, while locality makes the corresponding local scalar operators commute. The time-ordered product is consequently unchanged by the change of variables $x_{r}\mapsto\Lambda x_{r}$. Taking the endpoints to infinity removes their frame-dependent description, and each $S_{n}$ obeys
\[
  U(\Lambda)S_{n}U(\Lambda)^{-1}=S_{n}.
\]
The sum $S$ is Lorentz-invariant. The dimensions also match: $[\dd^{D}x]=-D$ and $[\mathcal L_{I}]=D$, so the exponent is dimensionless.

% BANKS-SOURCE: pdf=25 print=15 kind=solution id=2.8
\BanksSolution{2.8}
In $D=4$, let $k=(m,\mathbf0)$ and choose $L(p)$ with $L(p)k=p$, where $p^2=-m^2$. For a Lorentz transformation $\Lambda$, define
\[
  R_{W}(\Lambda,p)=L^{-1}(\Lambda p)\Lambda L(p).
\]
It leaves $k$ fixed:
\[
  R_{W}k=L^{-1}(\Lambda p)\Lambda p=L^{-1}(\Lambda p)(\Lambda p)=k.
\]
The stabilizer of a massive rest momentum is the rotation group, so $R_{W}$ is a rotation. Factoring $\Lambda L(p)=L(\Lambda p)R_{W}$ gives
\begin{align*}
  U(\Lambda)\ket{\mathbf p,k}
  &=\sqrt{\frac{m}{\omega_{\mathbf p}}}\,U(L(\Lambda p))U(R_{W})\ket{0,k}\\
  &=\sqrt{\frac{m}{\omega_{\mathbf p}}}\,D^{j}_{kl}(R_{W})U(L(\Lambda p))\ket{0,l}\\
  &=\sqrt{\frac{\omega_{\boldsymbol{\Lambda p}}}{\omega_{\mathbf p}}}
    D^{j}_{kl}(R_{W})\ket{\boldsymbol{\Lambda p},l}.
\end{align*}

If $a^{\dagger}_{k}(\mathbf p)\ket{0}=\ket{\mathbf p,k}$, the corresponding operator law is
\[
  U(\Lambda)a^{\dagger}_{k}(\mathbf p)U(\Lambda)^{-1}
  =\sqrt{\frac{\omega_{\boldsymbol{\Lambda p}}}{\omega_{\mathbf p}}}
    D^{j}_{kl}(R_{W})a^{\dagger}_{l}(\boldsymbol{\Lambda p}),
\]
with the annihilation law obtained by taking the adjoint. Let $S_{L}^{(j)}$ and $S_{R}^{(j)}$ denote the $(j,0)$ and $(0,j)$ Lorentz matrices. Rest-frame intertwiners $u(0)$ and $v(0)$ can be boosted to
\[
  u(p)=S_{L}^{(j)}(L(p))u(0),
  \qquad v(p)=S_{R}^{(j)}(L(p))v(0).
\]
They obey the intertwining relations needed to cancel the Wigner matrices in the operator law. A covariant field is therefore obtained from
\[
  \Psi_{L,A}(x)=\int\frac{\dd^{D-1}\mathbf p}{\sqrt{(2\pi)^{D-1}2\omega_{\mathbf p}}}
  \left[u_{A}{}^{k}(p)a_{k}(\mathbf p)\ee^{\ii p\mathbin{\cdot}x}
  +v_{A}{}^{k}(p)a^{\dagger}_{k}(\mathbf p)\ee^{-\ii p\mathbin{\cdot}x}\right],
\]
and its conjugate $\Psi_{R}=\Psi_{L}^{\dagger}$, which transform in $(j,0)$ and $(0,j)$, respectively. In spinor notation $A$ is a totally symmetric string of $2j$ undotted or dotted indices. Choosing the creation coefficients as the conjugates of the annihilation coefficients makes the doubled field Hermitian for a neutral particle.

The exchange sign is fixed by the little-group rotation carried by the intertwiners. Reversing two space-like separated field arguments produces the phase $\ee^{2\pi\ii j}=(-1)^{2j}$. If $\varepsilon$ is the sign in the graded locality relation, locality requires
\[
  \varepsilon=(-1)^{2j}.
\]
Integer spin therefore uses commuting oscillators, while half-integer spin uses anticommuting oscillators.

For an invariant adjoint $\overline\Psi$, the free two-point function has the form
\[
  \langle0|T\Psi_{A}(x)\overline\Psi_{B}(y)|0\rangle
  =\int\frac{\dd^{D}p}{(2\pi)^{D}}
  \frac{-\ii\,\mathcal N^{(j)}_{AB}(p)}{p^{2}+m^{2}-\ii\epsilon}
  \ee^{\ii p\mathbin{\cdot}(x-y)},
\]
where on shell $\mathcal N^{(j)}_{AB}(p)=\sum_{s}u_{A}(p,s)\overline u_{B}(p,s)$, with the appropriate graded time ordering. In the spinor cases, $\slashed p\equiv\gamma^{\mu}p_{\mu}$ and $\{\gamma^{\mu},\gamma^{\nu}\}=-2\eta^{\mu\nu}$ in $D=4$. Its highest-momentum term has degree $2j$ in $p$. Thus
\[
  \mathcal G^{(j)}(p)=O(|p|^{2j-2})\qquad(|p|\to\infty).
\]
For $j=0$ this is $p^{-2}$, for $j=\tfrac12$ it is $p^{-1}$, for $j=1$ it is $p^{0}$, and the powers rise to $p$ for $j=\tfrac32$ and $p^{2}$ for $j=2$. The familiar numerators $1$, $\slashed p+m$, and $\eta_{\mu\nu}+p_{\mu}p_{\nu}/m^{2}$ illustrate the first three cases.

% BANKS-SOURCE: pdf=25 print=15 kind=solution id=2.9
\BanksSolution{2.9}
An explicit matrix basis satisfying the stated condition is
\[
  (J_{\mu\nu})^{\rho}{}_{\sigma}
  =-\ii\left(\delta_{\mu}^{\rho}\eta_{\nu\sigma}
  -\delta_{\nu}^{\rho}\eta_{\mu\sigma}\right),
  \qquad J_{\mu\nu}=-J_{\nu\mu}.
\]
Indeed, direct multiplication gives $J\eta+\eta J^{T}=0$. Multiplying two such matrices and subtracting in the opposite order gives
\[
  [J_{\mu\nu},J_{\alpha\beta}]^{\rho}{}_{\sigma}
  =\ii\left(\eta_{\mu\alpha}(J_{\nu\beta})^{\rho}{}_{\sigma}
  -\eta_{\nu\alpha}(J_{\mu\beta})^{\rho}{}_{\sigma}
  +\eta_{\nu\beta}(J_{\mu\alpha})^{\rho}{}_{\sigma}
  -\eta_{\mu\beta}(J_{\nu\alpha})^{\rho}{}_{\sigma}\right),
\]
which is the displayed Lorentz algebra.

In $D=4$, set $K_{i}=J_{0i}$ and $J_{ij}=\epsilon_{ijk}J_{k}$. With $\eta_{00}=-1$ and $\eta_{ij}=\delta_{ij}$, the independent component relations are
\[
  [K_{i},K_{j}]=-\ii J_{ij}=-\ii\epsilon_{ijk}J_{k},
\]
\[
  [K_{i},J_{jk}]=-\ii(\delta_{ij}K_{k}-\delta_{ik}K_{j}),
\]
\[
  [J_{ij},J_{kl}]=\ii\left(\delta_{ik}J_{jl}-\delta_{jk}J_{il}
  +\delta_{jl}J_{ik}-\delta_{il}J_{jk}\right).
\]
For $\epsilon_{123}=+1$, the last two equations can also be written
\[
  [K_{i},J_{j}]=\ii\epsilon_{ijk}K_{k},
  \qquad
  [J_{i},J_{j}]=\ii\epsilon_{ijk}J_{k}.
\]
The signs in this final form follow the source convention $J_{ij}=+\epsilon_{ijk}J_{k}$ together with the mostly-plus metric.

% BANKS-SOURCE: pdf=25 print=15 kind=solution id=2.10
\BanksSolution{2.10}
In $D=4$, choose the standard null momentum $k=(\kappa,0,0,\kappa)$ with $\kappa>0$. In terms of $K_{i}=J_{0i}$ and the rotations $J_{i}$, a convenient basis for its Lorentz stabilizer is
\[
  J_{3},\qquad N_{1}=K_{1}+J_{2},\qquad N_{2}=K_{2}-J_{1}.
\]
Using the component algebra gives
\[
  [N_{1},N_{2}]=0,
  \qquad [J_{3},N_{1}]=\ii N_{2},
  \qquad [J_{3},N_{2}]=-\ii N_{1}.
\]
The generator $J_{3}$ rotates the pair $(N_{1},N_{2})$, and the two commuting generators give translations. These are the defining relations of $\mathfrak{iso}(2)$, so the connected little group is the Euclidean group $ISO(2)$.

In a finite-dimensional unitary representation, $N_{1}$ and $N_{2}$ are commuting Hermitian matrices. A nonzero joint eigenvalue would be rotated through a continuous circle by $J_{3}$, producing infinitely many distinct eigenvalues in a finite-dimensional space. Hence
\[
  N_{1}=N_{2}=0.
\]
The remaining generator has an eigenvalue $h$, the helicity. A $2\pi$ rotation acts as $1$ on single-valued bosonic states and as $-1$ on spinorial fermionic states. Therefore
\[
  h\in\mathbb Z\quad\text{for bosons},
  \qquad
  h\in\mathbb Z+\tfrac12\quad\text{for fermions}.
\]

Let $L(p)k=p$ for a positive-energy null momentum. The induced states
\[
  \ket{\mathbf p,h}=\sqrt{\frac{\kappa}{\omega_{\mathbf p}}}\,U(L(p))\ket{k,h}
\]
transform as
\[
  U(\Lambda)\ket{\mathbf p,h}
  =\sqrt{\frac{\omega_{\boldsymbol{\Lambda p}}}{\omega_{\mathbf p}}}
   \ee^{\ii h\vartheta(\Lambda,p)}\ket{\boldsymbol{\Lambda p},h},
\]
where $\vartheta(\Lambda,p)$ is the rotation angle of the $ISO(2)$ Wigner element
\[
  W(\Lambda,p)=L^{-1}(\Lambda p)\Lambda L(p).
\]
The translation part acts trivially. If $a_{h}^{\dagger}(\mathbf p)\ket0=\ket{\mathbf p,h}$, then
\[
  U(\Lambda)a_{h}^{\dagger}(\mathbf p)U(\Lambda)^{-1}
  =\sqrt{\frac{\omega_{\boldsymbol{\Lambda p}}}{\omega_{\mathbf p}}}
   \ee^{\ii h\vartheta(\Lambda,p)}a_{h}^{\dagger}(\boldsymbol{\Lambda p}),
\]
and $a_{h}(\mathbf p)$ carries the complex-conjugate phase.

A local finite-dimensional Lorentz field contains both weights $+h$ and $-h$. In two-component notation the corresponding smallest field strengths are totally symmetric spinors
\[
  \Phi_{\alpha_{1}\cdots\alpha_{2h}}\in(h,0),
  \qquad
  \overline\Phi_{\dot\alpha_{1}\cdots\dot\alpha_{2h}}\in(0,h),
\]
which carry the two helicities. For $h=1$ these are $F_{\alpha\beta}$ and $\overline F_{\dot\alpha\dot\beta}$, the self-dual and anti-self-dual parts of the electromagnetic field strength $F_{\mu\nu}$; their engineering dimension is two. For $h=\tfrac32$ they are the symmetric rank-three spinor curvatures of a free Rarita--Schwinger field, of dimension $\tfrac52$. For $h=2$ they are the symmetric rank-four spinors $C_{\alpha\beta\gamma\delta}$ and $\overline C_{\dot\alpha\dot\beta\dot\gamma\dot\delta}$, the self-dual and anti-self-dual Weyl curvatures, of dimension three. These local curvatures display the required $\pm h$ helicities and the representation content.

% BANKS-SOURCE: pdf=26 print=16 kind=solution id=2.11
\BanksSolution{2.11}
Split the free field into its annihilation and creation parts,
\[
\begin{aligned}
  \phi(x)&=\phi_{\rm a}(x)+\phi_{\rm c}(x),\\
  \phi_{\rm a}(x)&=\int\frac{\dd^{D-1}\mathbf p}{\sqrt{(2\pi)^{D-1}2\omega_{\mathbf p}}}
  a(\mathbf p)\ee^{\ii p\mathbin{\cdot}x},\\
  \phi_{\rm c}(x)&=\phi_{\rm a}(x)^{\dagger}.
\end{aligned}
\]
The only nonzero contraction is a c-number. With
\[
  \Delta(x-y)=\int\frac{\dd^{D}p}{(2\pi)^{D}}
  \frac{\ee^{\ii p\mathbin{\cdot}(x-y)}}{p^{2}+m^{2}-\ii\epsilon},
\]
the chapter-wide convention is
\[
  \langle0|T\phi(x)\phi(y)|0\rangle=-\ii\Delta(x-y).
\]
Wick's theorem then gives, for $A=\int\dd^{D}x\,J(x)\phi(x)$,
\[
  \langle0|T\ee^{\ii A}|0\rangle
  =\sum_{r=0}^{\infty}\frac{1}{r!}
    \left[-\frac12\int\dd^{D}x\,\dd^{D}y\,J(x)J(y)\,
    \langle0|T\phi(x)\phi(y)|0\rangle\right]^{r}.
\]
For example, the first terms are
\[
\begin{aligned}
  1&+\frac12\int\dd^{D}x\,\dd^{D}y\,J(x)J(y)\,\ii\Delta(x-y)\\
  &+\frac18\left(\int\dd^{D}x\,\dd^{D}y\,J(x)J(y)\,\ii\Delta(x-y)\right)^{2}+\cdots.
\end{aligned}
\]
All odd terms vanish because a vacuum matrix element with an odd number of free creation and annihilation operators has no complete pairing. The even terms have $(2r-1)!!$ pairings, which gives the factorial in the preceding series. Hence the operator calculation gives
\[
  \langle0|T\ee^{\ii\int J\phi}|0\rangle
  =\exp\!\left[\frac{\ii}{2}\int\dd^{D}x\,\dd^{D}y\,J(x)J(y)
  \int\frac{\dd^{D}p}{(2\pi)^{D}}
  \frac{\ee^{\ii p\mathbin{\cdot}(x-y)}}{p^{2}+m^{2}-\ii\epsilon}\right].
\]
In the source mostly-minus convention, the corrected exponent is $-\ii/2$ times the source kernel $1/(p_{\rm old}^{2}-m^{2}+\ii\epsilon)$. Since $p_{\rm old}^{2}=-p^{2}$, this becomes the displayed $+\ii/2$ with the adapted kernel $1/(p^{2}+m^{2}-\ii\epsilon)$.

For the static source, write $R^{\mu}=(0,\mathbf R)$ and use the time interval $-T<t<T$. At large $T$, the time integrations project onto zero energy. With
\[
\begin{aligned}
  G_{m,D}(R)&=\int\frac{\dd^{D-1}\mathbf q}{(2\pi)^{D-1}}
  \frac{\ee^{\ii\mathbf q\mathbin{\cdot}\mathbf R}}{\mathbf q^{2}+m^{2}}\\
  &=\frac{m^{(D-3)/2}}{(2\pi)^{(D-1)/2}R^{(D-3)/2}}
  K_{(D-3)/2}(mR),
  \qquad R=|\mathbf R|.
\end{aligned}
\]
the source Fourier factor is $|1-\ee^{\ii\mathbf q\mathbin{\cdot}\mathbf R}|^{2}=2(1-\cos(\mathbf q\mathbin{\cdot}\mathbf R))$. The quadratic exponent has an $R$-independent self-energy and an interaction term proportional to $G_{m,D}(R)$:
\[
  \log Z[J]=\text{self-energy}-2\ii T\,G_{m,D}(R)+o(T)
\]
for the corrected $+\ii/2$ convention. After absorbing the self-energy into the normalization, this gives
\[
  Z[J]=\ee^{-\ii T V(R)},
  \qquad
  V(R)=2G_{m,D}(R),
  \qquad V(R)\big|_{D=4}=\frac{\ee^{-mR}}{2\pi R}.
\]
Using the uncorrected source-convention sign would reverse the interaction term. In either convention, the $R$ dependence is the Yukawa form in $D=4$.

\end{bankssolutions}
